\subsection{Perfis de aproximação: R-BAR e V-BAR}

Os perfis de aproximação empregados em missões de RVD/B são as trajetórias ao longo da R-BAR (radial) e de V-BAR (\emph{along-track}).  
A aproximação pela radial se dá pelo movimento ao longo do eixo que conecta o centro da Terra ao alvo.  

\begin{figure}[ht]
\centering
\includegraphics[width=0.8\textwidth]{2 Formulacao matematica/Imagens/Aprox Vbar.png}
\caption{Perfil de aproximação tangencial (V-bar).}
\label{fig:aproximacao_vbar}
\end{figure}

A aproximação pelo \emph{along-track} envolve aproximação tangencial ao vetor de velocidade orbital, porém possui maior sensibilidade a perturbações orbitais.
Neste trabalho, adota-se a aproximação pelo eixo radial, devido a ser mais utilizada pelas missões comerciais atuais.

\begin{figure}[ht]
\centering
\includegraphics[width=0.8\textwidth]{2 Formulacao matematica/Imagens/Aprox Rbar.png}
\caption{Perfil de aproximação radial (R-bar).}
\label{fig:aproximacao_rbar}
\end{figure}